% !PRINTED 30
\weberpage{30}
But in order to prove that $R$ is a commutative group, we must first set out two auxiliary considerations.

1) It is to be proved that any two of the groups $Q_\nu, Q'_\nu, Q''_\nu \ldots$ have no common divisor other than the unit.

For this we need only prove that $Q_\nu, Q'_\nu$ have this property. For if
\[ Q'_\nu = \gamma'^{-1} Q_\nu \gamma', \quad Q''_\nu = \gamma''^{-1} Q_\nu \gamma'' \]
have the greatest common divisor $T$, then $Q_\nu$ and $Q'''_\nu = \gamma' Q''_\nu \gamma'^{-1}$ have the greatest common divisor $\gamma' T \gamma'^{-1}$.

If we assume, then, that $T$ is the greatest common divisor of $Q_\nu$ and $Q'_\nu$, then $T$ is certainly a commutative group, because $Q_\nu, Q'_\nu$ are such. But $T$ is also a normal divisor of $P_\nu$. For $Q_\nu$ and $Q'_\nu$ are such divisors; and if therefore $\tau$ is an element of $T$ and $\pi$ an element of $P_\nu$, then $\pi^{-1}\tau\pi$ is contained both in $Q_\nu$ and in $Q'_\nu$, and hence also in $T$. Consequently $\pi^{-1} T \pi = T$.

Now we have assumed, however, that $Q_\nu$ is a commutative normal divisor of $P_\nu$ above the unit of the lowest possible degree; further, $Q_\nu$ and $Q'_\nu$ were different from one another. Hence $T$ can only be the unit group itself, q. e. d.

2) To form the group $R$, we may proceed in this way: we compound the elements of $Q_\nu, Q'_\nu, Q''_\nu \ldots$ with one another, in arbitrary order and with repetition, so long as new elements still arise. For all the compositions so formed constitute a group $C$, which is contained in $R$, because $R$ contains the individual elements of $Q_\nu, Q'_\nu, Q''_\nu \ldots$. On the other hand, the groups $Q_\nu, Q'_\nu, Q''_\nu \ldots$ are also contained in $C$, and consequently $R$ is contained in $C$ and $C$ is identical with $R$.

In order, then, finally to prove that $R$ is a commutative group, it only remains for us to show that, if $\alpha, \beta$ are any two elements of $Q_\nu, Q'_\nu, Q''_\nu \ldots$, the permutability
\weq{10}{\alpha\beta = \beta\alpha}
holds.

Now if $\alpha, \beta$ both belong to the same group $Q_\nu$, then (10) is a consequence of our assumption that $Q_\nu$ is commutative; and it is likewise so if both occur in one of the other groups $Q'_\nu, Q''_\nu \ldots$.
